\section{Related Work}

Our 
work on scalable mastery learning ~\cite{Bloom68}, inspired by
Garcia et al.'s ``A's for All''~\cite{garcia-asforall} approach,
requires
the ability to administer frequent and
short summative
assessments that afford the
possibility of re-takes to demonstrate mastery on earlier
material~\cite{secondchanceexams}.
We believed that the extensive and
successful experience reported by the University of Illinois in
efficiently operating Computer-Based Testing
Facilities~\cite{zilles-testing-less-trying} was the most promising
approach to addressing that aspect of mastery learning.

Downey et al.~\cite{downey-cbtf} recently reported their experience
setting up a similar center at the University of California,
Riverside.
Our goals and experience are comparable to theirs, but whie they were
able to obtain full-time use of a secure room for their testing center
from the outset, we were limited to part-time use of existing
shared lab spaces.
In addition, their experience report was written after operating their
testing center for two academic quarters; the lessons reported here reflect two years of
hindsight and the experience of 20x scale-up, especially regarding
institution-wide concerns such as 
accessibility workflows, funding, and messaging/policy
issues.

Finally, mastery learning notwithstanding, many instructors who have
gravitated to our CBTF since the pilot were motivated 
by the need to secure assessments against unauthorized use
of AI.  In this context, the CBTF can be seen as the secure ``lane''
of the University of Sydney's ``two-lane approach to
assessment''~\cite{twolane}.
